\section{Conclusion}
A Hall-silent conducting sector can matter because its elimination changes both the relaxation and drive of Hall-active distributions. Active reference current and orthogonal valley distortion distinguish scalar rescaling from Hall-sensitive redistribution. Passive current and passive valley polarization transmit their kinetic influence; the latter is essential for the harder finite-range cases but decouples in the massless contact limit.

The specified contact kernel closes analytically on these physical coordinates, with only three driven amplitudes needed. Finite range breaks the factorization of incoming moments and introduces angular outgoing rates. Four modes remain accurate on the tested domain, but their finite leakage is a diagnostic whose observable effect also depends on drive miss, inverse relaxation and Hall overlap. The exact residual identity explains why a small Hall error can coexist with a larger distribution error.

Omission is robustly safe under the fixed-budget test when the absolute component budget cannot exceed the tolerance after favorable signs are removed. Otherwise a safe response can be cancellation-dependent under that test: 571/1410 (40.50\%), 377/1942 (19.41\%) and 136/2201 (6.18\%) of safe positive-coupling points at 1\%, 5\% and 10\%. Finite passive curvature adds a direct population readout that may dominate even before coupling, and conductivity changes the omission decision for open-circuit proxies. The resulting criterion is kinetic, observable specific and tolerance dependent within the anomalous-velocity model. The calculation establishes neither a complete experimental nonlinear Hall conductivity nor a material-independent omission rule.
